\section{Síntesis y ampliación}

\begin{enumerate}
    \item El diseño de recompensas es el cuello de botella de las aplicaciones de RL; reward learning ofrece una solución automatizada.
    \item \textbf{Goal Classifiers}: aprenden a partir de ejemplos de éxito y de fracaso; son simples y directos. El entrenamiento adversarial evita el reward hacking.
    \item \textbf{Inverse RL}: infiere la función de recompensa a partir del comportamiento experto; es más flexible que el aprendizaje por imitación.
    \item \textbf{Aprendizaje de preferencias}: entrena un modelo de recompensa a partir de comparaciones de preferencias humanas; es el núcleo de RLHF.
    \item Todos los métodos de reward learning enfrentan el riesgo de reward hacking, por lo que requieren regularización y restricciones.
\end{enumerate}

\subsection{Tabla comparativa de métodos}

\begin{table}[H]
\centering
\begin{tabular}{p{0.18\textwidth} p{0.25\textwidth} p{0.25\textwidth} p{0.22\textwidth}}
\toprule
Método & Señal central & Ventaja principal & Riesgo principal \\
\midrule
Goal Classifier & Probabilidad de clasificación binaria entre estados de éxito y de fracaso & Sencillo; fácil de aplicar con conjuntos de datos pequeños & Reward hacking, deriva de distribución \\
Inverse RL & Función de reward que optimizan las trayectorias expertas & Permite explicar la intención del experto & Requiere trayectorias de alta calidad; entrenamiento inestable \\
Aprendizaje de preferencias (RLHF) & winner vs loser elegidos por personas & Se alinea directamente con las preferencias humanas & Reward hacking, ruido en las evaluaciones \\
\bottomrule
\end{tabular}
\caption{Comparación de los métodos de reward learning}
\end{table}

\subsection{Lecturas complementarias y líneas de práctica}

\begin{itemize}
    \item Dar seguimiento a las memorias del reward modeling workshop del Stanford RL group
    \item Estudiar el protocolo de recolección de datos de «Deep RL from Human Preferences»
    \item Vincular el pipeline de registro propio con el reward trace, para establecer un control de versiones que permita revertir cambios
\end{itemize}

\subsection{Lecturas adicionales}
\begin{itemize}
    \item Christiano et al., \textit{Deep Reinforcement Learning from Human Preferences} (2017).
    \item Ziegler et al., \textit{Fine-Tuning Language Models from Human Preferences} (2019).
    \item Sharma et al., \textit{Self-Improving Robots: End-to-End Autonomous Visuomotor Reinforcement Learning} (2023).
\end{itemize}

\end{document}
